\section{Introduction}
\label{sec:introduction}

Robotic throwing extends the effective reach of a manipulator far beyond its kinematic
workspace, enabling rapid sorting, bin-to-bin transfer, and long-range placement without
additional actuators.
Industrial pick-and-throw cells can double throughput compared with pick-and-place
because the robot releases the object at high speed and immediately returns for the next
grasp while the projectile is still in flight~\cite{turcato2025mcpilot}.
Despite this appeal, practical deployment is bottlenecked by \emph{data efficiency}:
learning a reliable release policy on physical hardware must converge in tens of trials
rather than thousands, because each trial incurs wear, reconfiguration overhead, and
safety risk.
Standard model-free RL methods (\eg TossingBot~\cite{tossbot}) require orders of
magnitude more samples than this budget permits, motivating model-based approaches that
explicitly represent uncertainty and plan through it.

\mcpilot{}~\cite{turcato2025mcpilot} addresses this challenge by combining
\mcpilco{}~\cite{mcpilco}---a Gaussian-process (GP) dynamics model with Monte Carlo
policy gradients---with two throwing-specific modifications.
First, the policy is reduced to a single-shot speed map $\pi_\theta(\mathbf{P})$
that fires exactly once at $t{=}0$ given the target position $\mathbf{P}$, eliminating
the need for closed-loop feedback during flight.
Second, gripper release delay is estimated by Bayesian Optimisation, compensating for
systematic timing errors that would otherwise corrupt every trajectory.
The augmented state $\tilde{\mathbf{x}} = [\mathbf{x},\, \mathbf{P}]$ lets the GP learn
geometry-conditioned corrections, and the RBF policy network enables smooth interpolation
across the target workspace.
On a Franka Panda arm with a Gazebo/ROS simulator, \mcpilot{} reaches near-100\% target
coverage in approximately ten physical throws.

Despite these results, several practical gaps remain unexplored.
\textbf{(1) Geometry sensitivity.}
The published hyperparameters are calibrated to a single release height and target range;
na\"ive transfer to different heights silently degrades performance.
\textbf{(2) Arm portability.}
The paper is locked to a single robot platform; it is unknown whether the algorithm's
data-efficiency transfers to other manipulators.
\textbf{(3) Noise structure.}
Physical robots exhibit velocity slip, timing jitter, and gripper compliance whose
statistical structure differs from zero-mean additive noise; the learnability of each
noise type is not characterised.
\textbf{(4) Environmental disturbances.}
Outdoor or warehouse environments introduce wind, yet the algorithm has not been tested
under ambient airflow.
\textbf{(5) Perception.}
Deployment requires autonomous bin detection; the original work pre-specifies target
coordinates rather than detecting them visually.
\textbf{(6) Reproducibility.}
The original implementation requires Gazebo and ROS, limiting accessibility; the
algorithm's core physics is ballistic and admits a lightweight numpy simulation.

\paragraph{Contributions.}
\begin{enumerate}
  \item \textbf{Independent simulation reproduction.}
        A self-contained ballistic simulator and PyBullet visualiser that replicate
        \mcpilot{} without Gazebo or ROS, with all adaptation decisions documented.

  \item \textbf{Four diagnosed failure modes.}
        Systematic identification and correction of four non-obvious silent failure modes:
        broken gradient path, post-landing GP hallucination, cost-lengthscale saturation,
        and simulation-horizon unit mismatch.

  \item \textbf{Two convergence rules.}
        (a) The \emph{lengthscale-to-range scaling rule} $\ell_s \approx 0.15 \times
        \text{target\_range}$ prevents RBF blindness across all tested geometries.
        (b) The \emph{stratified exploration rule} guarantees seed-robust GP coverage
        with no additional hardware throws.

  \item \textbf{Release-height generalisation.}
        100\% hit rates across all five tested heights ($z \in \{0.0, 0.5, 1.0, 1.5,
        2.0\}$\,m) with mean landing errors below 35\,mm.

  \item \textbf{Multi-arm portability.}
        100\% commanded-velocity hit rates on Franka Panda, KUKA iiwa7, and xArm6
        without any retraining, demonstrating arm-agnostic RL through a decoupled
        architecture.

  \item \textbf{Zero-mean noise invariance theorem.}
        Theorem~\ref{thm:zeromean} proves that symmetric zero-mean noise does not shift
        the optimal policy, with an experimental demonstration of 100\% vs.\ 0\%
        hit rate for noise-aware vs.\ naive policies under 20\% velocity slip.

  \item \textbf{Wind robustness study.}
        A full pipeline extension (air-relative drag, 8-D$\to$11-D state augmentation,
        policy and GP updates) and the counterintuitive finding that blind GP policies
        outperform wind-aware policies in all three wind regimes within the 15-trial budget.

  \item \textbf{Closed-loop vision result.}
        We show that colour-based (OpenCV HSV) and learning-based (YOLOv8n on 800
        synthetic PyBullet images) detection pipelines both produce target estimates
        accurate enough to drive the trained policy to successful throws, with the
        near-face depth bias correctable by a geometric $2R/\pi$ shift.
        This establishes that the trained policy is robust to centimetre-level
        target-position estimation error.
\end{enumerate}

\paragraph{Paper structure.}
Section~\ref{sec:background} provides an expanded treatment of MC-PILCO and the
\mcpilot{} formulation.
Section~\ref{sec:system} describes the system architecture.
Section~\ref{sec:implementation} details the four diagnosed failure modes.
Section~\ref{sec:convergence} presents the two convergence rules.
Sections~\ref{sec:study1}--\ref{sec:study5} present the five experimental studies.
Section~\ref{sec:discussion} synthesises findings.
Section~\ref{sec:conclusion} concludes.
